\hypertarget{conclusion}{%
\section{Conclusion}\label{conclusion}}

The exponential cost of estimating \(\mathrm{Tr}(\rho^k)\) from local classical shadows is a
property of the estimator, and its rate is a property of the state. Both statements are made
precise by evaluating the two projection variances the Hoeffding decomposition already isolates:
\(\zeta_2\) is a Pauli sum whose kernel \(14^{-|P|}\) confines the per-qubit base to
\([7, 17/2]\), \(\zeta_1\) is a cubic contraction of the same second-moment identity, and the
budget threshold their ratio defines has base \(\mathrm{base}(\zeta_2)/\mathrm{base}(\zeta_1)\)
exactly, under a nonzero-leading-coefficient condition that the maximally mixed state violates.
Neither bound on \(\zeta_2\) is attained, and the finite-size ratio \(\zeta_2^{1/n}\) can sit
below \(7\); it is the limit that the bound governs.

Two consequences are worth separating from the machinery. Purity does not determine the rate: a
Haar-random pure state and a pure product state, both of purity one, sit on branches at \(28/5\)
and at exactly \(5\), and for every pure product state the threshold is available in closed form
at every \(n\). And the popular value \(28/5\) is the spread-spectrum case rather than a constant,
which is the sort of thing that is easy to over-generalize from one ensemble.

Because these are functionals of a state rather than of a family, they can be checked statewise,
and we did: across a family whose thresholds span a decade the law orders individual states with
rank correlation \(0.87\) and a slope of \(0.97\). The negative control matters as much as the
result. On the fixed-estimand families that the crossover literature, including our own earlier
sweeps, has used, the statewise spread is smaller than the achievable measurement noise, so no
such claim could have been established there --- a family-average law would have scored the same.

The threshold is also reachable from data. Unbiased estimators built from four disjoint blocks of
a pilot budget recover \(M^*\) without forming the Pauli spectrum, to ten percent from a pilot
that grows about twofold per qubit against the threshold's roughly fivefold, so the overhead
becomes relatively cheaper with size and falls below the threshold by eight qubits. What is
missing is a performance guarantee, which needs a concentration bound for a difference of degree-2
and degree-4 statistics and is a separate problem.

Paired with two exact bias laws for the collective route, distinguished by the geometry of the
noise channel rather than by its rate, the variance law predicts, with no parameter fitted to the
crossover outcomes, the size at which collective measurement attains the lower RMSE at equal copy
budget. It places every resolving cell within one qubit across five state families, two of which
were built for this paper because the other three cannot test a statewise claim. For estimating
purity of noisy-pure states under per-qubit amplitude-damping or dephasing noise at the rates and
budgets tested here, single-copy classical shadows stop being the lower-RMSE route beyond roughly
six to seven qubits; the higher moments hold out to about eight for the idealized bounded-outcome
test, whose \(k \ge 3\) realization costs about an order of magnitude more two-qubit gates
(Section~\ref{sec:compilation}), and under global depolarizing noise the crossover falls earlier
still.

Three questions follow. Whether the higher-order projection variances admit the same treatment,
which would extend the exact account past \(k = 2\) and is the obstacle to a performance
guarantee. Whether the same variance law, applied to the partial-transpose moments, gives an exact
sample-complexity account of entanglement negativity estimation, which is the application that
motivated us. And whether the branch structure --- base \(7\) for spread spectra, \(15/2\) for
stabilizer-structured ones --- has a classification behind it, since the two values we can prove
exactly are the two extremes of the families we happened to examine.
